In this section, we first present the results obtained for the \texttt{memcmps} code.\\
Next, we show the results for the \texttt{verify\_pin} code.\\
Finally, we discuss the findings related to the control-flow graph generated during the analysis.
%------------------------------------------------compare-----------------------------------------------
\subsection{memcmps}
\label{section:memcmps_result}
This section presents the results for the \texttt{memcmps} test.

\begin{table}[H]
    \centering
    \resizebox{0.4\textwidth}{!}{
        \begin{tabular}{l|l|l|l|l|l}
            \toprule
            Version & IP & Info                & 0-order & 1-order  & 2-order    \\
            \midrule
            V1      & 0  & data [l:19, bb:bb0] & 0       & 4 (mc:1) & 20 (mc:1)  \\
                    & 1  & ti [l:20, bb:bb0]   & 0       & 4 (mc:1) & 12 (mc:1)  \\
                    & 2  & data [l:21, bb:bb1] & 0       & 0        & 8 (mc:1)   \\
                    & 3  & data [l:23, bb:bb2] & 0       & 4 (mc:1) & 8 (mc:1)   \\
                    & 4  & data [l:24, bb:bb3] & 0       & 4 (mc:1) & 16 (mc:1)  \\
            \midrule
            V2      & 0  & data [l:7, bb:bb0]  & 0       & 0        & 24 (mc:1)  \\
                    & 1  & ti [l:7, bb:bb0]    & 0       & 0        & 8 (mc:1)   \\
                    & 2  & data [l:9, bb:bb1]  & 0       & 0        & 14 (mc:1)  \\
                    & 3  & data [l:10, bb:bb1] & 0       & 0        & 0          \\
                    & 4  & ti [l:10, bb:bb1]   & 0       & 0        & 0          \\
                    & 5  & data [l:12, bb:bb2] & 0       & 0        & 0          \\
                    & 6  & data [l:13, bb:bb2] & 0       & 0        & 0          \\
                    & 7  & ti [l:13, bb:bb2]   & 0       & 0        & 0          \\
                    & 8  & data [l:14, bb:bb3] & 0       & 0        & 0          \\
                    & 9  & data [l:19, bb:bb6] & 0       & 4 (mc:1) & 18 (mc:1)  \\
            \midrule
            V3      & 0  & data [l:7, bb:bb0]  & 0       & 0        & 24 (mc:1)  \\
                    & 1  & ti [l:7, bb:bb0]    & 0       & 0        & 8 (mc:1)   \\
                    & 10 & ti [l:13, bb:bb3]   & 0       & 0        & 0          \\
                    & 11 & data [l:14, bb:bb4] & 0       & 0        & 0          \\
                    & 12 & data [l:19, bb:bb8] & 0       & 4 (mc:1) & 18 (mc:1)  \\
                    & 2  & data [l:8, bb:bb1]  & 0       & 0        & 14 (mc:1)  \\
                    & 3  & data [l:9, bb:bb1]  & 0       & 0        & 0          \\
                    & 4  & ti [l:9, bb:bb1]    & 0       & 0        & 0          \\
                    & 5  & data [l:10, bb:bb2] & 0       & 0        & 0          \\
                    & 6  & data [l:11, bb:bb2] & 0       & 0        & 0          \\
                    & 7  & ti [l:11, bb:bb2]   & 0       & 0        & 0          \\
                    & 8  & data [l:12, bb:bb3] & 0       & 0        & 0          \\
                    & 9  & data [l:13, bb:bb3] & 0       & 0        & 0          \\
            \bottomrule
        \end{tabular}       
        
    }
    \captionof{table}{Table of attack for memcmps in C}
    \label{tab1:memcmps_C}
\end{table}

Table \ref{tab1:memcmps_C}: presents the attack found by Lazart for the code C.\\
Version 1 (V1) uses one mask and one call to \texttt{memcmp}.\\
Version 2 (V2) uses one mask and three calls to \texttt{memcmp}.\\
Version 3 (V3) uses two masks and four calls to \texttt{memcmp}.\\
Version 3 is also the one shown in Figure \ref{figure 1:memcmps_C}.\\
In the Info column: data (modification of loaded values) and ti (test inversion) indicate the type of attack used.
The attack are followed by the line in the code (l:int) and the block in control-flow graph (bb:name\_of\_block).\\
The Order column corresponds to the number of fault injections introduced in a single attack.

\begin{table}[H]
    \centering
    \resizebox{0.4\textwidth}{!}{
        \begin{tabular}{l|l|l|l|l|l}
            \toprule
            Version & IP & Info                & 0-order & 1-order  & 2-order    \\
            V1      & 0  & data [l:30, bb:start]       & 0       & 0        & 8 (mc:1)   \\
                    & 1  & ti [l:8, bb:bb2.i]          & 0       & 0        & 4 (mc:1)   \\
                    & 2  & ti [l:10, bb:bb4.i]         & 0       & 0        & 13 (mc:1)  \\
                    & 3  & data [l:30, bb:memcmp.exit] & 0       & 4 (mc:1) & 25 (mc:1)  \\
                    \midrule
            V2      & 0  & data [l:24, bb:start]       & 0       & 4 (mc:1) & 60 (mc:1)  \\
                    & 1  & ti [l:40, bb:bb2.i]         & 0       & 0        & 36 (mc:1)  \\
                    & 10 & ti [l:42, bb:bb4.i15]       & 0       & 0        & 0          \\
                    & 11 & ti [l:30, bb:memcmp.exit16] & 0       & 0        & 0          \\
                    & 2  & ti [l:42, bb:bb4.i]         & 0       & 0        & 17 (mc:1)  \\
                    & 3  & ti [l:24, bb:memcmp.exit]   & 0       & 0        & 8 (mc:1)   \\
                    & 4  & data [l:36, bb:bb11]        & 0       & 4 (mc:1) & 35 (mc:1)  \\
                    & 5  & data [l:26, bb:bb4]         & 0       & 0        & 22 (mc:1)  \\
                    & 6  & ti [l:40, bb:bb2.i2]        & 0       & 0        & 16 (mc:1)  \\
                    & 7  & ti [l:42, bb:bb4.i7]        & 0       & 0        & 0          \\
                    & 8  & ti [l:27, bb:memcmp.exit8]  & 0       & 0        & 0          \\
                    & 9  & ti [l:40, bb:bb2.i10]       & 0       & 0        & 16 (mc:1)  \\
                    \midrule
            V3      & 0  & data [l:24, bb:start]       & 0       & 4 (mc:1) & 76 (mc:1)  \\
                    & 1  & ti [l:41, bb:bb2.i]         & 0       & 0        & 36 (mc:1)  \\
                    & 10 & ti [l:43, bb:bb4.i15]       & 0       & 0        & 0          \\
                    & 11 & ti [l:28, bb:memcmp.exit16] & 0       & 0        & 0          \\
                    & 12 & ti [l:41, bb:bb2.i18]       & 0       & 0        & 16 (mc:1)  \\
                    & 13 & ti [l:43, bb:bb4.i23]       & 0       & 0        & 0          \\
                    & 14 & ti [l:30, bb:memcmp.exit24] & 0       & 0        & 0          \\
                    & 2  & ti [l:43, bb:bb4.i]         & 0       & 0        & 17 (mc:1)  \\
                    & 3  & ti [l:24, bb:memcmp.exit]   & 0       & 0        & 8 (mc:1)   \\
                    & 4  & data [l:37, bb:bb14]        & 0       & 4 (mc:1) & 35 (mc:1)  \\
                    & 5  & data [l:25, bb:bb4]         & 0       & 0        & 22 (mc:1)  \\
                    & 6  & ti [l:41, bb:bb2.i2]        & 0       & 0        & 16 (mc:1)  \\
                    & 7  & ti [l:43, bb:bb4.i7]        & 0       & 0        & 0          \\
                    & 8  & ti [l:26, bb:memcmp.exit8]  & 0       & 0        & 0          \\
                    & 9  & ti [l:41, bb:bb2.i10]       & 0       & 0        & 16 (mc:1)  \\
            \bottomrule
        \end{tabular}
    }
    \captionof{table}{Table of attack for memcmps in Rust}
    \label{tab2:memcmps_Rust}
    \end{table}

\noindent Table \ref{tab2:memcmps_Rust}: show us the attack found by lazart for the code Rust.\\

In Version 3,we observed that a single fault injection as create one attack at line 19 in Figure \ref{figure 1:memcmps_C} and at line 37 in Figure \ref{figure 2:memcmps_Rust}. Both attacks correspond to the return value of the function, which results in a mutation of the data stored in the \texttt{result} variable.
However, we also noted that Version 3 exhibits fewer attacks in Rust (see Table \ref{tab2:memcmps_Rust}) compared to C (see Table \ref{tab1:memcmps_C}).An explanation for this discrepancy will be provided in Control-Flow graph (\ref{subsection:result cfg}) and Conclusion (\ref{section:Conclusion}) section.

%------------------------------------------------verify_pin-----------------------------------------------
\subsection{Verify pin}
This section presents the results for the \texttt{verifypin} test.
\begin{table}[ht]
    \centering
    \resizebox{0.4\textwidth}{!}{
        \begin{tabular}{l|l|l|l|l|l} 
            \toprule
            Version & IP & Info                                       & 0-order & 1-order   & 2-order     \\ 
            \midrule
            0       & 0  & ti [l:8, bb:bb0]                           & 0       & 4 (mc:1)  & 6 (mc:1)    \\
                    & 1  & ti [l:10, bb:check in compare]             & 0       & 4 (mc:1)  & 27 (mc:2)   \\
                    & 2  & ti [l:20, bb:Try counter \textgreater{} 0] & 0       & 0         & 0           \\
                    & 3  & ti [l:23, bb:precompare]                   & 0       & 4 (mc:1)  & 9 (mc:1)    \\ 
            \midrule
            1       & 0  & ti [l:8, bb:bb0]                           & 0       & 4 (mc:1)  & 6 (mc:1)    \\
                    & 1  & ti [l:9, bb:bb1]                           & 0       & 4 (mc:1)  & 27 (mc:2)   \\
                    & 2  & ti [l:20, bb:Try counter \textgreater{} 0] & 0       & 0         & 0           \\
                    & 3  & ti [l:24, bb:precompare]                   & 0       & 4 (mc:1)  & 9 (mc:1)    \\
                    & 4  & ti [l:27, bb:bb7]                          & 0       & 0         & 0           \\ 
            \midrule
            2       & 0  & ti [l:10, bb:bb0]                          & 0       & 0         & 4 (mc:1)    \\
                    & 1  & ti [l:13, bb:check in compare]             & 0       & 4 (mc:1)  & 124 (mc:2)  \\
                    & 2  & ti [l:19, bb:i!=size]                      & 0       & 0         & 4 (mc:1)    \\
                    & 3  & ti [l:24, bb:diff in compare]              & 0       & 15 (mc:1) & 56 (mc:1)   \\
                    & 4  & ti [l:37, bb:Try counter \textgreater{} 0] & 0       & 0         & 0           \\
                    & 5  & ti [l:41, bb:precompare]                   & 0       & 15 (mc:1) & 56 (mc:1)   \\
                    & 6  & ti [l:45, bb:bb8]                          & 0       & 0         & 0           \\ 
            \midrule
            4       & 0  & ti [l:12, bb:Try counter \textgreater{} 0] & 0       & 0         & 0           \\
                    & 1  & ti [l:13, bb:bb0]                          & 0       & 0         & 0           \\
                    & 10 & ti [l:53, bb:bb13]                         & 0       & 0         & 0           \\
                    & 2  & ti [l:18, bb:bb1]                          & 0       & 0         & 0           \\
                    & 3  & ti [l:29, bb:bb2]                          & 0       & 0         & 0           \\
                    & 4  & ti [l:30, bb:bb3]                          & 0       & 4 (mc:1)  & 124 (mc:2)  \\
                    & 5  & ti [l:35, bb:bb6]                          & 0       & 0         & 0           \\
                    & 6  & ti [l:38, bb:bb7]                          & 0       & 0         & 0           \\
                    & 7  & ti [l:41, bb:bb8]                          & 0       & 15 (mc:1) & 56 (mc:1)   \\
                    & 8  & ti [l:47, bb:bb10]                         & 0       & 15 (mc:1) & 56 (mc:1)   \\
                    & 9  & ti [l:48, bb:bb11]                         & 0       & 0         & 0           \\ 
            \midrule
            6       & 0  & ti [l:12, bb:Try counter \textgreater{} 0] & 0       & 0         & 0           \\
                    & 1  & ti [l:19, bb:bb0]                          & 0       & 0         & 1 (mc:1)    \\
                    & 2  & ti [l:22, bb:user\_pin != card\_pin]       & 0       & 0         & 0           \\
                    & 3  & ti [l:28, bb:i != PIN\_SIZE]               & 0       & 0         & 1 (mc:1)    \\
                    & 4  & ti [l:33, bb:diff 1]                       & 0       & 0         & 1 (mc:1)    \\
                    & 5  & ti [l:36, bb:diff 2]                       & 0       & 0         & 1 (mc:1)    \\
                    & 6  & ti [l:48, bb:status 1]                     & 0       & 0         & 1 (mc:1)    \\
                    & 7  & ti [l:51, bb:status 2]                     & 0       & 0         & 1 (mc:1)    \\
            \bottomrule
        \end{tabular}
    }
    \captionof{table}{Table of attack for verifypin in C}
    \label{tab3:verifypin_C}
\end{table}
\noindent Table \ref{tab3:verifypin_C}: present the attacks identified by Lazart for the code C.

This table \ref{tab3:verifypin_C} is similar to the one discribed in the previous subsection \ref{section:memcmps_result}.
But it corresponds to a different set of countermeasures implemented in \texttt{verify\_pin} test:
\begin{itemize}
        \item \textbf{Version 0} corresponds to \texttt{verify\_pin} which does not include any countermesure.
        \item \textbf{Version 1} uses only hardened bool values.
        \item \textbf{Version 2} combine a fixed-time loop with hardened boolean value.
        \item \textbf{Version 4} introduced :
        \begin{itemize}
                \item An anticipated decrement of the \texttt{`try\_counter`} variable (try\_counter is performed in an separate conditional block).
                \item A duplication of the \texttt{`try\_counter`} variable.
                \item The two additional countermesures from Version 2.
        \end{itemize}
        \item \textbf{Version 6} combines :
        \begin{itemize}
                \item An anticipated decrement of the \texttt{`try\_counter`} variable.
                \item A fixed-time loop.
                \item Hardened boolean values.
                \item Test duplication (the PIN is tested twice in two different blocks).
        \end{itemize} 
\end{itemize}

%\onecolumn
%\begin{multicols}{2}
 
\begin{table}[H]
    \centering
    \resizebox{0.4\textwidth}{!}{
        \begin{tabular}{l|l|l|l|l|l} 
            \toprule
            Version & IP & Info                                       & 0-order & 1-order  & 2-order    \\ 
            \midrule
            0       & 0  & ti [l:8, bb:for\_in\_compare]              & 0       & 1 (mc:1) & 0          \\
                    & 1  & ti [l:515, bb:check in compare]            & 0       & 4 (mc:1) & 26 (mc:2)  \\
                    & 2  & ti [l:8, bb:\_ZN70\_\$LT\$core..option…]   & 0       & 3 (mc:1) & 14 (mc:1)  \\
                    & 3  & ti [l:21, bb:Try counter \textless{} 0]    & 0       & 0        & 0          \\
                    & 4  & ti [l:24, bb:precompare]                   & 0       & 1 (mc:1) & 8 (mc:1)   \\ 
            \midrule
            1       & 0  & ti [l:8, bb:for\_in\_compare]              & 0       & 1 (mc:1) & 0          \\
                    & 1  & ti [l:515, bb:check in compare]            & 0       & 0        & 0          \\
                    & 2  & ti [l:515, bb:bb2.i]                       & 0       & 4 (mc:1) & 22 (mc:2)  \\
                    & 3  & ti [l:8, bb:\_ZN70\_\$LT\$core..option…]   & 0       & 3 (mc:1) & 10 (mc:1)  \\
                    & 4  & ti [l:22, bb:Try counter \textless{} 0]    & 0       & 0        & 0          \\ 
            \midrule
            2       & 0  & ti [l:10, bb:for\_in\_compare]             & 0       & 0        & 1 (mc:1)   \\
                    & 1  & ti [l:17, bb:i!=size]                      & 0       & 0        & 27 (mc:1)  \\
                    & 2  & ti [l:515, bb:check in compare]            & 0       & 0        & 0          \\
                    & 3  & ti [l:515, bb:bb2.i]                       & 0       & 0        & 4 (mc:1)   \\
                    & 4  & ti [l:12, bb:bb6.i]                        & 0       & 0        & 4 (mc:1)   \\
                    & 5  & ti [l:10, bb:bb13]                         & 0       & 0        & 3 (mc:1)   \\
                    & 6  & ti [l:35, bb:Try counter \textless{} 0]    & 0       & 0        & 0          \\
                    & 7  & ti [l:38, bb:precompare]                   & 0       & 0        & 15 (mc:1)  \\
                    & 8  & ti [l:43, bb:bb9]                          & 0       & 0        & 0          \\ 
            \midrule
            4       & 0  & ti [l:13, bb:Try counter \textgreater{} 0] & 0       & 0        & 0          \\
                    & 1  & ti [l:34, bb:for loop]                     & 0       & 0        & 1 (mc:1)   \\
                    & 2  & ti [l:34, bb:bb23]                         & 0       & 0        & 1 (mc:1)   \\
                    & 3  & ti [l:34, bb:bb23.1]                       & 0       & 0        & 1 (mc:1)   \\
                    & 4  & ti [l:34, bb:bb23.2]                       & 0       & 0        & 1 (mc:1)   \\
                    & 5  & ti [l:42, bb:bb23.3]                       & 0       & 0        & 0          \\
                    & 6  & ti [l:47, bb:bb28]                         & 0       & 0        & 34 (mc:1)  \\
                    & 7  & ti [l:52, bb:bb33]                         & 0       & 0        & 15 (mc:1)  \\
                    & 8  & ti [l:61, bb:bb39]                         & 0       & 0        & 15 (mc:1)  \\
                    & 9  & ti [l:70, bb:bb47]                         & 0       & 0        & 0          \\ 
            \midrule
            6       & 0  & ti [l:10, bb:Try counter \textgreater{} 0] & 0       & 0        & 0          \\
                    & 1  & ti [l:18, bb:for loop]                     & 0       & 0        & 1 (mc:1)   \\
                    & 2  & ti [l:28, bb:i != PIN\_SIZE]               & 0       & 0        & 12 (mc:1)  \\
                    & 3  & ti [l:515, bb:user\_pin != card\_pin]      & 0       & 0        & 4 (mc:1)   \\
                    & 4  & ti [l:21, bb:bb6.i]                        & 0       & 0        & 4 (mc:1)   \\
                    & 5  & ti [l:18, bb:bb16]                         & 0       & 0        & 3 (mc:1)   \\
                    & 6  & ti [l:34, bb:diff 1]                       & 0       & 0        & 0          \\
                    & 7  & ti [l:36, bb:diff 2]                       & 0       & 0        & 0          \\
                    & 8  & ti [l:48, bb:status 1]                     & 0       & 0        & 0          \\
                    & 9  & ti [l:51, bb:status 2]                     & 0       & 0        & 0          \\
            \bottomrule
        \end{tabular}
    }
    \captionof{table}{Table of attack for verifypin in Rust}
    \label{tab4:verifypin_Rust}
\end{table}
\noindent Table \ref{tab4:verifypin_Rust}: present the attacks identified by Lazart for the Rust code.\\

We observed similar results with \texttt{verify\_pin} test in both Table \ref{tab3:verifypin_C} and Table \ref{tab4:verifypin_Rust}.

%------------------------------------------------CFG-----------------------------------------------

       
\subsection{Control Flow Graph}
\label{subsection:result cfg}
We also observed some major differences in how Rust compiles compared to C.\\
The Intermediate code generated by LLVM revealed that Rust optimises certain conditional blocks differently than C.\\

However, we noticed differences in the control-flow graphs generated by the two language versions of the code presented in this study.\\
Specifically, Rust transforms variables with two possible values—depending on the evaluation of a condition, into a \texttt{select} instruction rather than a branch instruction.
Similarly, Rust may use a \texttt{switch} instruction when a variable is compared to multiple values.